\DocumentMetadata{
  pdfversion=2.0,
  pdfstandard=ua-2,
  lang=en-US,
  tagging=on,
  tagging-setup={math/setup=mathml-SE,tabsorder=structure}
}
\documentclass[11pt,letterpaper]{article}
\usepackage[margin=0.68in,includefoot]{geometry}
\usepackage{newcomputermodern}
\usepackage{amsmath,amscd,mathtools}
\usepackage{tikz,adjustbox}
\providecommand{\square}{\mdlgwhtsquare}
\usepackage[hidelinks]{hyperref}
\hypersetup{
  pdftitle={Analysis PhD Exam, January 2019},
  pdfauthor={University of Florida Department of Mathematics},
  pdfsubject={Graduate mathematics examination},
  pdfdisplaydoctitle=true,
  bookmarksopen=true
}
\setcounter{secnumdepth}{0}
\setlength{\parindent}{0pt}
\setlength{\parskip}{0.28em}
\setlength{\emergencystretch}{3em}
\setlength{\leftmargini}{2.15em}
\setlength{\leftmarginii}{2.65em}
\setlength{\leftmarginiii}{3em}
\pagestyle{empty}
\raggedbottom
\allowdisplaybreaks
\NewTaggingSocketPlug{block/list/label}{examproblem}{%
  \tagstructbegin{tag=Lbl}\tagstructend%
  \tagstructbegin{tag=\LBody}%
  \tagstructbegin{tag=\ProblemHeadingTag,title-o={\ProblemHeadingText}}%
  \tagmcbegin{tag=\ProblemHeadingTag,actualtext={\ProblemHeadingText}}%
  #2%
  \tagmcend%
  \tagstructend%
}
\newenvironment{examproblems}[2]{%
  \def\ProblemHeadingTag{#2}%
  \AssignTaggingSocketPlug{block/list/label}{examproblem}%
  \begin{enumerate}%
  \setcounter{enumi}{#1}%
  \renewcommand{\labelenumi}{\textbf{\theenumi.}}%
  \setlength{\topsep}{0.35em}%
  \setlength{\partopsep}{0pt}%
  \setlength{\itemsep}{0.48em}%
  \setlength{\parsep}{0.18em}%
}{\end{enumerate}}
\newenvironment{examsubparts}{%
  \AssignTaggingSocketPlug{block/list/label}{default}%
  \begin{enumerate}%
  \setlength{\topsep}{0.18em}%
  \setlength{\partopsep}{0pt}%
  \setlength{\itemsep}{0.16em}%
  \setlength{\parsep}{0.1em}%
}{\end{enumerate}}
\newenvironment{examinstructions}{%
  \par\begingroup\itshape
}{%
  \par\endgroup\vspace{0.18em}
}
\makeatletter
\renewcommand\subsection{\@startsection{subsection}{2}{0pt}%
  {-1.25ex plus -.25ex minus -.1ex}%
  {0.55ex plus .1ex}%
  {\normalfont\large\bfseries}}
\renewcommand{\@maketitle}{%
  \newpage\null\vskip -1em%
  {\footnotesize\bfseries
    \href{https://www.ufl.edu/}{University of Florida}, \href{https://math.ufl.edu/}{Department of Mathematics}\par}%
  \vskip -0.5em%
  \tagstructbegin{tag=H1,title={Analysis PhD Exam, January 2019}}%
  \tagpdfparaOff%
  \tagmcbegin{tag=H1}%
    {\LARGE\bfseries\href{https://gma.math.ufl.edu/exams/analysis-phd/}{Analysis PhD Exam}, January 2019\par}%
  \tagmcend%
  \tagpdfparaOn%
  \tagstructend%
  \vskip 0.25em%
  \tagmcbegin{artifact}%
    \hrule height 1pt%
  \tagmcend%
  \vskip 1em%
}
\makeatother
\title{Analysis PhD Exam, January 2019}
\author{}
\date{}
\begin{document}
\maketitle
\thispagestyle{empty}
\begin{examinstructions}
Do three from part A and three from part B. Write solutions in a neat and logical fashion, giving complete reasons for all steps.
\end{examinstructions}
\subsection{Part A}
\begin{examproblems}{0}{H3}
\def\ProblemHeadingText{Problem 1.}
\item
Let \(X,Y\) be topological spaces. Prove that every continuous function \(f:X\to Y\) is Borel measurable.
\def\ProblemHeadingText{Problem 2.}
\item
Short answer.
\begin{examsubparts}
\item[\textnormal{(a)}]
Give an example, if possible, of a closed set \(E\subset\mathbb R\) of positive Lebesgue measure that contains no interval.
\item[\textnormal{(b)}]
For \(n\geq 2\), let \(f_n=e^{inx}x^{-n}\). Find
\[
\lim_{n\to\infty}\int_1^\infty f_n\,dx.
\]
\item[\textnormal{(c)}]
Give an example, if possible, of strict inequality in Fatou’s Lemma.
\end{examsubparts}
\def\ProblemHeadingText{Problem 3.}
\item
Let \(X=Y=[0,1]\), \(\mathcal M=\mathcal B_{[0,1]}\) denote the Borel~\(\sigma\)-algebra, and \(\mathcal N=2^{\mathbb R}\). Let~\(\mu\) Lebesgue measure on \(\mathcal M\), and let~\(\nu\) counting measure on \(\mathcal N\). Prove that \(\Delta=\{(x,x):x\in[0,1]\}\subset[0,1]\times[0,1]\) is an element of the product~\(\sigma\)-algebra \(\mathcal M\otimes\mathcal N\) and compute
\[
\int_X\left(\int_Y \mathbf 1_\Delta(x,y)\,d\nu(y)\right)d\mu(x),\qquad
\int_Y\left(\int_X \mathbf 1_\Delta(x,y)\,d\mu(x)\right)d\nu(y). \tag{1}
\]
 Explain the relation to Tonelli’s Theorem on product integration.
\def\ProblemHeadingText{Problem 4.}
\item
State the Hahn Decomposition Theorem. Prove, if~\(\rho\) is a signed measure on the measurable space \((X,\mathcal M)\), then there exist unique positive measures \(\rho_\pm\) on \(\mathcal M\) such that \(\rho_+\perp\rho_-\) and \(\rho=\rho_+-\rho_-\).

\par
Assuming~\(\mu\) is a positive measure on \(\mathcal M\) and \(f\in L^1(\mu)\) is real-valued, identify \(\rho_\pm\) for the signed measure
\[
\rho(E)=\int_E f\,d\mu.
\]
\end{examproblems}
\subsection{Part B}
\begin{examproblems}{4}{H3}
\def\ProblemHeadingText{Problem 5.}
\item
State the Baire Category Theorem. Suppose~\(X\) is a Banach space. Prove that, as a vector space, \(X\) does not have a countable basis. (Here countable means equinumerous with \(\mathbb N\).)
\def\ProblemHeadingText{Problem 6.}
\item
Let \((\varphi_n)\) be a sequence from \(L^1(\mathbb R)\) with the following properties:

\par
Show, if \(f\in L^1(\mathbb R)\), then \(\varphi_n\star f\) converges to~\(f\) in \(L^1(\mathbb R)\).
\begin{examsubparts}
\item[\textnormal{(i)}]
\(\varphi_n\geq 0\) for all~\(n\),
\item[\textnormal{(ii)}]
\(\displaystyle \int\varphi_n\,dx=1\) for all~\(n\);
\item[\textnormal{(iii)}]
for each \(\delta>0\), \(\displaystyle \lim_{n\to\infty}\int_{|x|>\delta}\varphi_n\,dx=0\).
\end{examsubparts}
\def\ProblemHeadingText{Problem 7.}
\item
Let~\(X\) be a Banach space, \(Y\subset X\) a closed subspace. Say~\(Y\) is complemented in~\(X\) if there exists a closed subspace \(Z\subset X\) such that \(Y\cap Z=\{0\}\) and \(Y+Z=X\).

\par
Prove that if~\(Y\) is complemented in~\(X\), then there exists a bounded linear operator \(T:X\to Y\) such that \(T(y)=y\) for all \(y\in Y\).
\def\ProblemHeadingText{Problem 8.}
\item
Let \(S:\ell^2(\mathbb N)\to\ell^2(\mathbb N)\) denote the shift operator defined, for \(f=(f(n))_{n=0}^\infty\in\ell^2(\mathbb N)\), by
\[
Sf(n)=
\begin{cases}
f(n-1),&n\geq 1;\\
0,&n=0.
\end{cases}
\]
 Show~\(S\) is an isometry (and in particular is bounded) and find \(S^*\) and \(SS^*\). Is~\(S\) unitary?
\end{examproblems}
\end{document}
